% =========================================
% frontmatter.tex — Declaration, Acks, Abstract, TOC/LoF/LoT
% Call from main.tex AFTER the title page:
%   % =========================================
% frontmatter.tex — Declaration, Acks, Abstract, TOC/LoF/LoT
% Call from main.tex AFTER the title page:
%   % =========================================
% frontmatter.tex — Declaration, Acks, Abstract, TOC/LoF/LoT
% Call from main.tex AFTER the title page:
%   % =========================================
% frontmatter.tex — Declaration, Acks, Abstract, TOC/LoF/LoT
% Call from main.tex AFTER the title page:
%   \input{frontmatter}
% =========================================

% Roman numbering for front matter
\pagenumbering{roman}

% -------- Declaration --------
\cleardoublepage
\addcontentsline{toc}{chapter}{Declaration}
\thispagestyle{empty}

\begin{center}
  {\Large \textbf{DECLARATION}}
\end{center}
\vspace{0.8cm}

I, \textbf{\AuthorName} (Roll No.\ \textbf{\RollNo}), hereby declare that the work presented
in the report entitled \textbf{``\ThesisTitle''}, submitted in partial fulfilment of the
requirements for the award of the degree of \textbf{Master of Technology (\Programme)} to
\textbf{\InstituteName}, is an authentic record of my own work carried out during
\textbf{\SessionDates} under the supervision of \textbf{\Supervisor} in the \textbf{\Department}.

\vspace{0.6cm}
This report has not been submitted for the award of any other degree or diploma of this or any
other institute. All sources of information used in the report—texts, figures, tables, and data—
have been duly acknowledged and cited.

\vspace{1.2cm}
\noindent
\begin{minipage}[t]{0.48\textwidth}
  \textbf{Date:} \rule{5cm}{0.4pt}\\[0.6cm]
  \textbf{Place:} \InstituteCity
\end{minipage}\hfill
\begin{minipage}[t]{0.48\textwidth}
  \raggedleft
  \textbf{\AuthorName}\\
  Roll No.: \textbf{\RollNo}
\end{minipage}

% -------- Acknowledgements (airy layout) --------
\cleardoublepage
\addcontentsline{toc}{chapter}{Acknowledgements}
\thispagestyle{empty}

\begin{center}
  {\Large \textbf{ACKNOWLEDGEMENTS}}
\end{center}
\vspace{0.9cm}

\begin{center}
\begin{minipage}{0.92\textwidth}
\begin{spacing}{1.2}
{\large
I express my deepest gratitude to my supervisor, \textbf{\Supervisor}, for
his patient guidance, insightful discussions, and continuous encouragement
throughout the course of this work. His clarity of thought and emphasis on
first principles shaped both the problem formulation and the execution of this
report.

I am grateful to the faculty and staff of the \textbf{\Department} at
\textbf{\InstituteName} for creating a supportive academic environment and for
their timely help with laboratory facilities, software access, and logistics.
I also thank the postgraduate office and administrative staff for their kind
assistance with official procedures.

My sincere thanks to my peers and labmates for the many helpful conversations,
code reviews, and constructive feedback that improved the methodology and
experiments presented here. I also acknowledge the broader open-source
communities around \textit{RISC\mbox{-}V}, toolchains, simulators, and FPGA
flows whose resources were invaluable.

Above all, I am profoundly thankful to my family and friends for their
unwavering support, patience, and motivation during the preparation of this
report.
}
\end{spacing}
\end{minipage}
\end{center}

\vspace{1.6cm}
\noindent
\begin{minipage}[t]{0.48\textwidth}
  \textbf{Date:} \rule{5cm}{0.4pt}\\[0.6cm]
  \textbf{Place:} \InstituteCity
\end{minipage}\hfill
\begin{minipage}[t]{0.48\textwidth}
  \raggedleft
  \textbf{\AuthorName}\\
  Roll No.: \textbf{\RollNo}
\end{minipage}

% -------- Abstract (compact, airy layout) --------
\cleardoublepage
\addcontentsline{toc}{chapter}{Abstract}
\thispagestyle{empty}

\begin{center}
  {\Large \textbf{ABSTRACT}}
\end{center}
\vspace{0.6cm}

\begin{center}
\begin{minipage}{0.92\textwidth}
\begin{spacing}{1.2}
{\large
This thesis presents a dual-architecture hardware implementation study of the Advanced
Encryption Standard (AES) on FPGA, targeting reproducible area and power baselines for
AES-128, AES-192, and AES-256. Two architectures are implemented and characterized: a
\textbf{full-pipeline} design that instantiates one copy of the round logic per pipeline
stage, and a \textbf{finite state machine (FSM)} design that reuses a single round hardware
block across all rounds under explicit state control. The pipeline architecture establishes a
high-throughput, high-area reference point (AES-128 baseline: 38,766 LUTs), while the FSM
architecture demonstrates area efficiency through hardware reuse (AES-128 baseline: 6,794
LUTs), representing an 82\% reduction.

Both architectures are subjected to two targeted optimization phases. \textbf{Phase~1}
replaces the EXP3 and LN3 log-table GF($2^{8}$) multiplier in the MixColumns datapath
with \textit{xtime} shift-and-XOR arithmetic, eliminating 512 lookup-table entries from
every MixColumns instance. \textbf{Phase~2} introduces clock enable guards on the FSM
state registers, suppressing unnecessary switching across all fields of the packed register
struct when the FSM is idle. Together, the two phases provide a controlled,
reproducible study of arithmetic simplification and power-aware gating across a
representative pair of FPGA AES architectures.
}
\end{spacing}
\end{minipage}
\end{center}

\vspace{0.8cm}
\noindent\textbf{Keywords---} AES; FPGA; pipeline architecture; FSM architecture;
GF($2^{8}$) arithmetic; xtime; clock enable gating; LUT utilization; hardware optimization;
AES-128; AES-192; AES-256.

% -------- TOC + Lists --------
\cleardoublepage
\tableofcontents

\cleardoublepage
\listoffigures

\cleardoublepage
\listoftables

% Switch to Arabic numbering for main matter
\cleardoublepage
\pagenumbering{arabic}

% =========================================

% Roman numbering for front matter
\pagenumbering{roman}

% -------- Declaration --------
\cleardoublepage
\addcontentsline{toc}{chapter}{Declaration}
\thispagestyle{empty}

\begin{center}
  {\Large \textbf{DECLARATION}}
\end{center}
\vspace{0.8cm}

I, \textbf{\AuthorName} (Roll No.\ \textbf{\RollNo}), hereby declare that the work presented
in the report entitled \textbf{``\ThesisTitle''}, submitted in partial fulfilment of the
requirements for the award of the degree of \textbf{Master of Technology (\Programme)} to
\textbf{\InstituteName}, is an authentic record of my own work carried out during
\textbf{\SessionDates} under the supervision of \textbf{\Supervisor} in the \textbf{\Department}.

\vspace{0.6cm}
This report has not been submitted for the award of any other degree or diploma of this or any
other institute. All sources of information used in the report—texts, figures, tables, and data—
have been duly acknowledged and cited.

\vspace{1.2cm}
\noindent
\begin{minipage}[t]{0.48\textwidth}
  \textbf{Date:} \rule{5cm}{0.4pt}\\[0.6cm]
  \textbf{Place:} \InstituteCity
\end{minipage}\hfill
\begin{minipage}[t]{0.48\textwidth}
  \raggedleft
  \textbf{\AuthorName}\\
  Roll No.: \textbf{\RollNo}
\end{minipage}

% -------- Acknowledgements (airy layout) --------
\cleardoublepage
\addcontentsline{toc}{chapter}{Acknowledgements}
\thispagestyle{empty}

\begin{center}
  {\Large \textbf{ACKNOWLEDGEMENTS}}
\end{center}
\vspace{0.9cm}

\begin{center}
\begin{minipage}{0.92\textwidth}
\begin{spacing}{1.2}
{\large
I express my deepest gratitude to my supervisor, \textbf{\Supervisor}, for
his patient guidance, insightful discussions, and continuous encouragement
throughout the course of this work. His clarity of thought and emphasis on
first principles shaped both the problem formulation and the execution of this
report.

I am grateful to the faculty and staff of the \textbf{\Department} at
\textbf{\InstituteName} for creating a supportive academic environment and for
their timely help with laboratory facilities, software access, and logistics.
I also thank the postgraduate office and administrative staff for their kind
assistance with official procedures.

My sincere thanks to my peers and labmates for the many helpful conversations,
code reviews, and constructive feedback that improved the methodology and
experiments presented here. I also acknowledge the broader open-source
communities around \textit{RISC\mbox{-}V}, toolchains, simulators, and FPGA
flows whose resources were invaluable.

Above all, I am profoundly thankful to my family and friends for their
unwavering support, patience, and motivation during the preparation of this
report.
}
\end{spacing}
\end{minipage}
\end{center}

\vspace{1.6cm}
\noindent
\begin{minipage}[t]{0.48\textwidth}
  \textbf{Date:} \rule{5cm}{0.4pt}\\[0.6cm]
  \textbf{Place:} \InstituteCity
\end{minipage}\hfill
\begin{minipage}[t]{0.48\textwidth}
  \raggedleft
  \textbf{\AuthorName}\\
  Roll No.: \textbf{\RollNo}
\end{minipage}

% -------- Abstract (compact, airy layout) --------
\cleardoublepage
\addcontentsline{toc}{chapter}{Abstract}
\thispagestyle{empty}

\begin{center}
  {\Large \textbf{ABSTRACT}}
\end{center}
\vspace{0.6cm}

\begin{center}
\begin{minipage}{0.92\textwidth}
\begin{spacing}{1.2}
{\large
This thesis presents a dual-architecture hardware implementation study of the Advanced
Encryption Standard (AES) on FPGA, targeting reproducible area and power baselines for
AES-128, AES-192, and AES-256. Two architectures are implemented and characterized: a
\textbf{full-pipeline} design that instantiates one copy of the round logic per pipeline
stage, and a \textbf{finite state machine (FSM)} design that reuses a single round hardware
block across all rounds under explicit state control. The pipeline architecture establishes a
high-throughput, high-area reference point (AES-128 baseline: 38,766 LUTs), while the FSM
architecture demonstrates area efficiency through hardware reuse (AES-128 baseline: 6,794
LUTs), representing an 82\% reduction.

Both architectures are subjected to two targeted optimization phases. \textbf{Phase~1}
replaces the EXP3 and LN3 log-table GF($2^{8}$) multiplier in the MixColumns datapath
with \textit{xtime} shift-and-XOR arithmetic, eliminating 512 lookup-table entries from
every MixColumns instance. \textbf{Phase~2} introduces clock enable guards on the FSM
state registers, suppressing unnecessary switching across all fields of the packed register
struct when the FSM is idle. Together, the two phases provide a controlled,
reproducible study of arithmetic simplification and power-aware gating across a
representative pair of FPGA AES architectures.
}
\end{spacing}
\end{minipage}
\end{center}

\vspace{0.8cm}
\noindent\textbf{Keywords---} AES; FPGA; pipeline architecture; FSM architecture;
GF($2^{8}$) arithmetic; xtime; clock enable gating; LUT utilization; hardware optimization;
AES-128; AES-192; AES-256.

% -------- TOC + Lists --------
\cleardoublepage
\tableofcontents

\cleardoublepage
\listoffigures

\cleardoublepage
\listoftables

% Switch to Arabic numbering for main matter
\cleardoublepage
\pagenumbering{arabic}

% =========================================

% Roman numbering for front matter
\pagenumbering{roman}

% -------- Declaration --------
\cleardoublepage
\addcontentsline{toc}{chapter}{Declaration}
\thispagestyle{empty}

\begin{center}
  {\Large \textbf{DECLARATION}}
\end{center}
\vspace{0.8cm}

I, \textbf{\AuthorName} (Roll No.\ \textbf{\RollNo}), hereby declare that the work presented
in the report entitled \textbf{``\ThesisTitle''}, submitted in partial fulfilment of the
requirements for the award of the degree of \textbf{Master of Technology (\Programme)} to
\textbf{\InstituteName}, is an authentic record of my own work carried out during
\textbf{\SessionDates} under the supervision of \textbf{\Supervisor} in the \textbf{\Department}.

\vspace{0.6cm}
This report has not been submitted for the award of any other degree or diploma of this or any
other institute. All sources of information used in the report—texts, figures, tables, and data—
have been duly acknowledged and cited.

\vspace{1.2cm}
\noindent
\begin{minipage}[t]{0.48\textwidth}
  \textbf{Date:} \rule{5cm}{0.4pt}\\[0.6cm]
  \textbf{Place:} \InstituteCity
\end{minipage}\hfill
\begin{minipage}[t]{0.48\textwidth}
  \raggedleft
  \textbf{\AuthorName}\\
  Roll No.: \textbf{\RollNo}
\end{minipage}

% -------- Acknowledgements (airy layout) --------
\cleardoublepage
\addcontentsline{toc}{chapter}{Acknowledgements}
\thispagestyle{empty}

\begin{center}
  {\Large \textbf{ACKNOWLEDGEMENTS}}
\end{center}
\vspace{0.9cm}

\begin{center}
\begin{minipage}{0.92\textwidth}
\begin{spacing}{1.2}
{\large
I express my deepest gratitude to my supervisor, \textbf{\Supervisor}, for
his patient guidance, insightful discussions, and continuous encouragement
throughout the course of this work. His clarity of thought and emphasis on
first principles shaped both the problem formulation and the execution of this
report.

I am grateful to the faculty and staff of the \textbf{\Department} at
\textbf{\InstituteName} for creating a supportive academic environment and for
their timely help with laboratory facilities, software access, and logistics.
I also thank the postgraduate office and administrative staff for their kind
assistance with official procedures.

My sincere thanks to my peers and labmates for the many helpful conversations,
code reviews, and constructive feedback that improved the methodology and
experiments presented here. I also acknowledge the broader open-source
communities around \textit{RISC\mbox{-}V}, toolchains, simulators, and FPGA
flows whose resources were invaluable.

Above all, I am profoundly thankful to my family and friends for their
unwavering support, patience, and motivation during the preparation of this
report.
}
\end{spacing}
\end{minipage}
\end{center}

\vspace{1.6cm}
\noindent
\begin{minipage}[t]{0.48\textwidth}
  \textbf{Date:} \rule{5cm}{0.4pt}\\[0.6cm]
  \textbf{Place:} \InstituteCity
\end{minipage}\hfill
\begin{minipage}[t]{0.48\textwidth}
  \raggedleft
  \textbf{\AuthorName}\\
  Roll No.: \textbf{\RollNo}
\end{minipage}

% -------- Abstract (compact, airy layout) --------
\cleardoublepage
\addcontentsline{toc}{chapter}{Abstract}
\thispagestyle{empty}

\begin{center}
  {\Large \textbf{ABSTRACT}}
\end{center}
\vspace{0.6cm}

\begin{center}
\begin{minipage}{0.92\textwidth}
\begin{spacing}{1.2}
{\large
This thesis presents a dual-architecture hardware implementation study of the Advanced
Encryption Standard (AES) on FPGA, targeting reproducible area and power baselines for
AES-128, AES-192, and AES-256. Two architectures are implemented and characterized: a
\textbf{full-pipeline} design that instantiates one copy of the round logic per pipeline
stage, and a \textbf{finite state machine (FSM)} design that reuses a single round hardware
block across all rounds under explicit state control. The pipeline architecture establishes a
high-throughput, high-area reference point (AES-128 baseline: 38,766 LUTs), while the FSM
architecture demonstrates area efficiency through hardware reuse (AES-128 baseline: 6,794
LUTs), representing an 82\% reduction.

Both architectures are subjected to two targeted optimization phases. \textbf{Phase~1}
replaces the EXP3 and LN3 log-table GF($2^{8}$) multiplier in the MixColumns datapath
with \textit{xtime} shift-and-XOR arithmetic, eliminating 512 lookup-table entries from
every MixColumns instance. \textbf{Phase~2} introduces clock enable guards on the FSM
state registers, suppressing unnecessary switching across all fields of the packed register
struct when the FSM is idle. Together, the two phases provide a controlled,
reproducible study of arithmetic simplification and power-aware gating across a
representative pair of FPGA AES architectures.
}
\end{spacing}
\end{minipage}
\end{center}

\vspace{0.8cm}
\noindent\textbf{Keywords---} AES; FPGA; pipeline architecture; FSM architecture;
GF($2^{8}$) arithmetic; xtime; clock enable gating; LUT utilization; hardware optimization;
AES-128; AES-192; AES-256.

% -------- TOC + Lists --------
\cleardoublepage
\tableofcontents

\cleardoublepage
\listoffigures

\cleardoublepage
\listoftables

% Switch to Arabic numbering for main matter
\cleardoublepage
\pagenumbering{arabic}

% =========================================

% Roman numbering for front matter
\pagenumbering{roman}

% -------- Declaration --------
\cleardoublepage
\addcontentsline{toc}{chapter}{Declaration}
\thispagestyle{empty}

\begin{center}
  {\Large \textbf{DECLARATION}}
\end{center}
\vspace{0.8cm}

I, \textbf{\AuthorName} (Roll No.\ \textbf{\RollNo}), hereby declare that the work presented
in the report entitled \textbf{``\ThesisTitle''}, submitted in partial fulfilment of the
requirements for the award of the degree of \textbf{Master of Technology (\Programme)} to
\textbf{\InstituteName}, is an authentic record of my own work carried out during
\textbf{\SessionDates} under the supervision of \textbf{\Supervisor} in the \textbf{\Department}.

\vspace{0.6cm}
This report has not been submitted for the award of any other degree or diploma of this or any
other institute. All sources of information used in the report—texts, figures, tables, and data—
have been duly acknowledged and cited.

\vspace{1.2cm}
\noindent
\begin{minipage}[t]{0.48\textwidth}
  \textbf{Date:} \rule{5cm}{0.4pt}\\[0.6cm]
  \textbf{Place:} \InstituteCity
\end{minipage}\hfill
\begin{minipage}[t]{0.48\textwidth}
  \raggedleft
  \textbf{\AuthorName}\\
  Roll No.: \textbf{\RollNo}
\end{minipage}

% -------- Acknowledgements (airy layout) --------
\cleardoublepage
\addcontentsline{toc}{chapter}{Acknowledgements}
\thispagestyle{empty}

\begin{center}
  {\Large \textbf{ACKNOWLEDGEMENTS}}
\end{center}
\vspace{0.9cm}

\begin{center}
\begin{minipage}{0.92\textwidth}
\begin{spacing}{1.2}
{\large
I express my deepest gratitude to my supervisor, \textbf{\Supervisor}, for
his patient guidance, insightful discussions, and continuous encouragement
throughout the course of this work. His clarity of thought and emphasis on
first principles shaped both the problem formulation and the execution of this
report.

I am grateful to the faculty and staff of the \textbf{\Department} at
\textbf{\InstituteName} for creating a supportive academic environment and for
their timely help with laboratory facilities, software access, and logistics.
I also thank the postgraduate office and administrative staff for their kind
assistance with official procedures.

My sincere thanks to my peers and labmates for the many helpful conversations,
code reviews, and constructive feedback that improved the methodology and
experiments presented here. I also acknowledge the broader open-source
communities around \textit{RISC\mbox{-}V}, toolchains, simulators, and FPGA
flows whose resources were invaluable.

Above all, I am profoundly thankful to my family and friends for their
unwavering support, patience, and motivation during the preparation of this
report.
}
\end{spacing}
\end{minipage}
\end{center}

\vspace{1.6cm}
\noindent
\begin{minipage}[t]{0.48\textwidth}
  \textbf{Date:} \rule{5cm}{0.4pt}\\[0.6cm]
  \textbf{Place:} \InstituteCity
\end{minipage}\hfill
\begin{minipage}[t]{0.48\textwidth}
  \raggedleft
  \textbf{\AuthorName}\\
  Roll No.: \textbf{\RollNo}
\end{minipage}

% -------- Abstract (compact, airy layout) --------
\cleardoublepage
\addcontentsline{toc}{chapter}{Abstract}
\thispagestyle{empty}

\begin{center}
  {\Large \textbf{ABSTRACT}}
\end{center}
\vspace{0.6cm}

\begin{center}
\begin{minipage}{0.92\textwidth}
\begin{spacing}{1.2}
{\large
This thesis presents a dual-architecture hardware implementation study of the Advanced
Encryption Standard (AES) on FPGA, targeting reproducible area and power baselines for
AES-128, AES-192, and AES-256. Two architectures are implemented and characterized: a
\textbf{full-pipeline} design that instantiates one copy of the round logic per pipeline
stage, and a \textbf{finite state machine (FSM)} design that reuses a single round hardware
block across all rounds under explicit state control. The pipeline architecture establishes a
high-throughput, high-area reference point (AES-128 baseline: 38,766 LUTs), while the FSM
architecture demonstrates area efficiency through hardware reuse (AES-128 baseline: 6,794
LUTs), representing an 82\% reduction.

Both architectures are subjected to two targeted optimization phases. \textbf{Phase~1}
replaces the EXP3 and LN3 log-table GF($2^{8}$) multiplier in the MixColumns datapath
with \textit{xtime} shift-and-XOR arithmetic, eliminating 512 lookup-table entries from
every MixColumns instance. \textbf{Phase~2} introduces clock enable guards on the FSM
state registers, suppressing unnecessary switching across all fields of the packed register
struct when the FSM is idle. Together, the two phases provide a controlled,
reproducible study of arithmetic simplification and power-aware gating across a
representative pair of FPGA AES architectures.
}
\end{spacing}
\end{minipage}
\end{center}

\vspace{0.8cm}
\noindent\textbf{Keywords---} AES; FPGA; pipeline architecture; FSM architecture;
GF($2^{8}$) arithmetic; xtime; clock enable gating; LUT utilization; hardware optimization;
AES-128; AES-192; AES-256.

% -------- TOC + Lists --------
\cleardoublepage
\tableofcontents

\cleardoublepage
\listoffigures

\cleardoublepage
\listoftables

% Switch to Arabic numbering for main matter
\cleardoublepage
\pagenumbering{arabic}
